\documentclass[../../main.tex]{subfiles}

\begin{document}

\section{参考文献}\label{section:ch12-参考文献}

D. Tataru 为本章提供了意见和建议。

\textbf{\S12.2}：Shatah--Struwe~\cite{S-S} 是本节的主要来源。

\textbf{\S12.3}：更多内容见 Shatah--Struwe~\cite{S-S}。\cref{theorem:ch12-符号条件全局存在} 遵循 I. Segal (Bull. Soc. Math. France 91 (1963), 129--135)；\cref{theorem:ch12-三维短时间爆破} 和\cref{theorem:ch12-次临界存在} 遵循 K. J\"orgens (Math. Z. 77 (1961), 295--308)。

\textbf{\S12.4}：M. Grillakis (Ann. of Math. 132 (1990), 485--509) 证明了\cref{theorem:ch12-L6能量通量估计} 中的关键估计和存在性\cref{theorem:ch12-u5全局存在}。我们对后者的介绍使用了 Smith 和 Sogge 的简化：参见 Sogge~\cite{So} 和 Shatah--Struwe~\cite{S-S}。

\textbf{\S12.5}：见 Strauss~\cite{St1}，John~\cite{J3}。

\textbf{\S12.6}：\cref{exercise:ch12-二维爆破} 来自 Protter--Weinberger~\cite{P-W}。\cref{exercise:ch12-小能量p5}：J. Rauch 观察到小能量蕴含 $p=5$ 时的全局存在性。J. Holmer 帮助完成了\cref{exercise:ch12-NLSGalilean}--\cref{exercise:ch12-NLS爆破}。

\end{document}